\icode{N20-EMI-SI-L3-261004-002}
\nceeID{2020全国\Rmnum{3}}{unkown}
如图，水平放置的圆形光滑玻璃棒左侧绕有一线圈，右侧套有一金属圆环，圆环初始时静止。将图中开关由断开状态拨至闭合状态，电路接通的瞬间，可观察到\paren[]

\begin{tikzpicture}[font=\scriptsize]
	\draw (-.2,0) rectangle(3.5,.5);
	\foreach \x in{.2,.6,...,3}
		\draw[blue] (\x,0)..controls({\x+.1},-.4)and({\x+.2},.9)..({\x+.3},.5);
	\draw (0,0.5)..controls(.05,.6)..(.1,.5);
	\draw[blue] (0,.5)--(0,-1.4)--(1.7,-1.4) (2.9,0)--(2.9,-2)--(1.5,-2)--(1.5,-1.6)node[black,right]{\(N\)} (2.9,-.8)--(1.5,-.8)--(1.5,-1.2)node[black,right]{\(M\)};
	\draw (3.3,.5) arc(40:325:.12 and .4);
	\draw[fill=white,draw=white] (2.2,-1)coordinate(battery) rectangle ($(battery)+(.1,.4)$);
    \draw[blue] ($(battery)+(0,.4)$)--++(0,-.4) ($(battery)+(.1,.3)$)--++(0,-.2);
	\draw[fill=white,draw=white] (2.2,-2.2)coordinate(battery) rectangle ($(battery)+(.1,.4)$);
    \draw[blue] ($(battery)+(0,.3)$)--++(0,-.2) ($(battery)+(.1,.4)$)--++(0,-.4);
	\foreach \c in{(1.2,-1.4),(1.5,-1.2),(1.5,-1.6)}
		\draw[draw=blue,fill=white] \c circle(.04);
	\node at(1.2,-1.4)[below]{\(S\)};
\end{tikzpicture}	
\begin{choices}
	\item 拨至\(M\)端或\(N\)端，圆环均向左运动
	\item 拨至\(M\)端或\(N\)端，圆环均向右运动
	\item 拨至\(M\)端，圆环向左运动；拨至\(N\)端，圆环向右运动
	\item 拨至\(M\)端，圆环向右运动；拨至\(N\)端，圆环向左运动
\end{choices}
\begin{answer}
    A\qquad 提示：无论拨至哪一端，圆环内的磁通量均增大，都要远离
\end{answer}